% Standalone insertable LaTeX appendix section for the FDFL-mu /
% PCGrad-normalization empirical follow-up. Intended to be \input{}'d
% into the IJOC manuscript. Assumes the manuscript already loads
% graphicx and booktabs.
% Expected manuscript-side figure path:
%   Plots/fig_gradient_scale_mu_sweep.png

\section{Empirical Follow-Up: Fixed Prediction Weights and Normalized PCGrad}
\label{sec:exp-mu-followup}

\subsection{Motivation}
The main experiments revealed a sharp contrast between the healthcare
and knapsack settings. On healthcare, where decision gradients are
computed analytically, end-to-end methods were consistently stable and
the scalarized and multi-objective variants clustered tightly. On the
synthetic multi-dimensional knapsack task, by contrast, the fair
decision-focused methods were much more sensitive to the
decision-gradient approximation. Per-iteration gradient logs showed two
distinct failure modes specific to the SPSA decision-gradient regime.
First, FDFL with no prediction weight ($\mu=0$) drifted during training
whenever the SPSA estimate was noisy; without an MSE anchor, the
predictor had no corrective signal to counteract stochastic decision
updates. Second, PCGrad saw per-objective gradient norms differing by
roughly three orders of magnitude---decision norms near $10^3$ and
prediction / fairness norms near $10^0$---so the projection geometry
that PCGrad uses to resolve conflicts was effectively dominated by a
single objective on every step.

Two targeted fixes follow directly from this diagnosis. The first is a
fixed prediction weight for FDFL, parameterized by $\mu$: we replace
the original loss $L_{\mathrm{regret}} + \lambda F$ with
\begin{equation}
  L_{\mathrm{FDFL}}(\mu) = L_{\mathrm{regret}}
  + \mu\, L_{\mathrm{pred}} + \lambda F.
  \label{eq:fdfl-mu}
\end{equation}
The original FDFL corresponds to $\mu=0$, and FDFL-Scal corresponds to
$\mu=1$; the intermediate values $\mu\in\{0.1,0.5\}$ allow us to test
whether even a small anchor is sufficient. The second fix is
per-objective normalization inside PCGrad: before the pairwise
conflict-projection loop, each objective gradient is rescaled to unit
$L_2$ norm, and the summed direction is then multiplied back by the
mean of the original norms so the overall step size is preserved. This
prevents any single objective's raw magnitude from biasing the
projection while leaving PCGrad's optimization step scale unchanged.

\subsection{Follow-up design}
We ran two follow-up experiments. The first is an MD-knapsack
``$\mu$-sweep'' on the redesigned per-individual, two-resource knapsack
with group-level $\alpha$-fair welfare, using $\alpha=2$, MAD
fairness, $n_{\mathrm{train}}=300$, 30 training steps per $\lambda$,
and SPSA decision gradients with 8 directions and perturbation scale
$5\times 10^{-3}$. The initial run used five seeds
$\{11,22,33,44,55\}$; a follow-up run on four additional seeds
$\{66,77,88,99\}$ was used to estimate the catastrophic-failure rate
of normalized PCGrad, yielding nine seeds in total. Ten methods were
evaluated: FPTO, SAA, WDRO, FDFL with $\mu\in\{0,0.1,0.5\}$,
FDFL-Scal, FPLG, PCGrad (with per-objective normalization), and MGDA.
The second experiment is a healthcare append run on the full Obermeyer
cohort with analytic decision gradients, which extends the existing
grid with the three fixed-$\mu$ FDFL variants while leaving the
paper-cited \texttt{stage\_results.csv} files untouched. The append
run spans four prediction-fairness measures, two welfare curvatures
$\alpha\in\{0.5,2\}$, five seeds, and 70 training steps per $\lambda$.
In both experiments, performance is summarized by normalized test
regret at the best $\lambda$ per (method, seed).

\subsection{Results on MD knapsack (SPSA)}
The $\mu$-sweep isolates the no-anchor variant as the unstable one.
Table~\ref{tab:mu-followup-mdknapsack} reports normalized test regret
for the main seed set at $\lambda=0$. FDFL with $\mu=0$ reaches
$0.838 \pm 0.356$, diverging on two of five seeds and attaining
genuine regret collapse (normalized regret $\approx 1.0$) in those
runs. Introducing even a small fixed prediction weight is sufficient
to restore stability: FDFL-$0.1$ falls to $0.197 \pm 0.004$ and
FDFL-$0.5$ to $0.186 \pm 0.004$, placing both variants in the same
performance band as FDFL-Scal ($0.175 \pm 0.007$) and FPLG
($0.191 \pm 0.004$). This confirms the diagnosis: under SPSA, the
decision gradient alone is too noisy to keep the predictor calibrated,
whereas any positive MSE weight prevents the runaway trajectory
observed for $\mu=0$.

Normalized PCGrad, in contrast, does not cleanly resolve the
SPSA-induced failure. On the main five seeds its summary statistic is
$0.354 \pm 0.362$, well above the anchored cluster and with very large
variance. The source of the variance is a single catastrophic seed
(seed~55) rather than uniformly noisy performance. To quantify this
effect we ran PCGrad on four additional seeds with the same task and
training configuration; all four produced stable runs (normalized
regret in $[0.163,\,0.207]$). Over the full nine-seed panel, eight
seeds are stable at $0.189 \pm 0.024$---numerically indistinguishable
from the anchored methods---while seed 55 alone reaches $0.9999$ with
an empirical training regret of $1.65\times 10^9$, a deterministic
runaway rather than gradient overflow. This indicates a residual
failure mode that normalization does not eliminate: the rescaling of
the summed direction by the mean of the original gradient norms still
allows an SPSA magnitude spike to inflate the effective step size, and
a single large step into a degenerate region of parameter space is
sufficient to trap the predictor, since subsequent SPSA estimates
around that region are equally uninformative. The observed failure
rate is $1/9 \approx 11\%$, not an artifact of the original five-seed
sample.

\begin{table}[t]
\centering
\small
\caption{MD-knapsack follow-up at $\lambda=0$. Columns report
normalized test regret as mean $\pm$ standard deviation over seeds.
FDFL with $\mu=0$ diverges in $2/5$ seeds; normalized PCGrad fails in
$1/9$ seeds, with the remaining eight seeds stable and competitive.
Lower is better. Bold marks the best anchored method per column.}
\label{tab:mu-followup-mdknapsack}
\begin{tabular}{lcc}
\toprule
Method & 5-seed mean & 9-seed / stable-subset mean \\
\midrule
FDFL ($\mu=0$)     & 0.838 $\pm$ 0.356     & --- \\
FDFL-0.1           & 0.197 $\pm$ 0.004     & --- \\
FDFL-0.5           & 0.186 $\pm$ 0.004     & --- \\
FDFL-Scal ($\mu=1$)& 0.175 $\pm$ 0.007     & --- \\
FPLG               & \textbf{0.191 $\pm$ 0.004} & --- \\
PCGrad (normalized)& 0.354 $\pm$ 0.362     & 0.189 $\pm$ 0.024 (8/9 stable) \\
MGDA               & 0.167 $\pm$ 0.006     & --- \\
\bottomrule
\end{tabular}
\end{table}

\subsection{Results on healthcare (analytic gradients)}
The healthcare append experiment tells a complementary story.
Table~\ref{tab:mu-followup-healthcare} reports normalized test regret
at the best $\lambda$ per (method, seed) for each
(fairness type, $\alpha$) cell. Two observations stand out.

First, the fixed-$\mu$ FDFL variants are numerically indistinguishable
from one another on this task. At MAD fairness and $\alpha=2$, FDFL
with $\mu=0$ attains $0.128\pm 0.002$, FDFL-$0.1$ attains
$0.128\pm 0.002$, FDFL-$0.5$ attains $0.129\pm 0.002$, and FDFL-Scal
attains $0.131\pm 0.002$. The same clustering holds across the other
fairness measures and at $\alpha=0.5$ (not shown individually, since
the within-cell mu-cluster spread never exceeds $0.003$). This is the
expected pattern: when the decision gradient is analytic and
low-variance, the predictor is pulled toward the correct outcomes
directly by the regret term, and the MSE anchor becomes redundant.

Second, normalized PCGrad has a strikingly $\alpha$-dependent effect.
At $\alpha=0.5$, normalization drives PCGrad well below the anchored
scalarized family across all four fairness measures, with relative
reductions in the 16--34\% range (for example, at MAD $\alpha=0.5$:
PCGrad $0.048\pm 0.001$ vs. FDFL-Scal $0.064$ and FPLG $0.065$). At
$\alpha=2$, the effect reverses: normalized PCGrad is a modest
$3$--$5$\% worse than the anchored cluster (for example, at MAD
$\alpha=2$: PCGrad $0.133\pm 0.004$ vs.\ FDFL-Scal $0.131$ and FPLG
$0.128$). The sign of the effect is consistent across all four
fairness types and both welfare curvatures, which suggests it is not
sampling noise. The most plausible explanation is that per-objective
normalization selectively advantages the concave-welfare regime, where
the prediction and fairness gradients are comparable in magnitude and
actively conflict, over the convex-welfare regime, where the
decision-regret gradient is already well aligned with the dominant
objective and equalizing magnitudes removes useful scale information.

\begin{table}[t]
\centering
\small
\caption{Healthcare follow-up, normalized test regret (mean $\pm$ std
over 5 seeds, best-$\lambda$ per seed). Across the four prediction
fairness measures, the FDFL $\mu$-family is indistinguishable within
each cell. Normalized PCGrad improves on the anchored cluster at
$\alpha=0.5$ across all four fairness measures and is marginally
worse at $\alpha=2$. Lower is better. Bold marks the best method per
column.}
\label{tab:mu-followup-healthcare}
\begin{tabular}{lcccccccc}
\toprule
& \multicolumn{2}{c}{MAD}
& \multicolumn{2}{c}{DP}
& \multicolumn{2}{c}{Atkinson}
& \multicolumn{2}{c}{Bias-parity} \\
\cmidrule(lr){2-3}\cmidrule(lr){4-5}\cmidrule(lr){6-7}\cmidrule(lr){8-9}
$\alpha$ & 0.5 & 2 & 0.5 & 2 & 0.5 & 2 & 0.5 & 2 \\
\midrule
FDFL           & 0.064 & 0.128 & 0.075 & 0.131 & 0.075 & 0.131 & 0.074 & 0.130 \\
FDFL-0.1       & 0.064 & 0.128 & 0.074 & 0.131 & 0.074 & 0.131 & 0.074 & 0.130 \\
FDFL-0.5       & 0.064 & 0.129 & 0.074 & 0.131 & 0.074 & 0.131 & 0.074 & 0.130 \\
FDFL-Scal      & 0.064 & 0.131 & 0.075 & 0.131 & 0.075 & 0.131 & 0.075 & 0.131 \\
FPLG           & 0.065 & \textbf{0.128} & 0.074 & \textbf{0.127} & 0.074 & \textbf{0.127} & 0.073 & \textbf{0.127} \\
PCGrad (norm.) & \textbf{0.048} & 0.133 & \textbf{0.063} & 0.144 & \textbf{0.053} & 0.137 & \textbf{0.049} & 0.130 \\
\bottomrule
\end{tabular}
\end{table}

\subsection{Gradient-scale diagnostic}
Figure~\ref{fig:gradient-scale-mu-followup} visualizes the mechanism
behind the contrasting outcomes. The top row (MD knapsack, SPSA)
shows a large separation between decision-gradient and
prediction-gradient norms, with a pronounced vertical blow-up of the
decision norm at $\mu=0$ and a sharp scale contraction for $\mu>0$.
The bottom row (healthcare, analytic gradients) is visually flat:
decision and prediction norms remain on the same order of magnitude
across all four $\mu$ values throughout training, and the trajectories
for $\mu\in\{0,0.1,0.5,1\}$ are nearly indistinguishable. The figure
reinforces the main message: a fixed prediction weight matters when
the decision gradient is noisy and poorly scaled, and has negligible
effect when the gradient is analytic and well-behaved.

\begin{figure}[t]
  \centering
  \includegraphics[width=0.92\textwidth]{Plots/fig_gradient_scale_mu_sweep.png}
  \caption{Gradient-scale diagnostic for the $\mu$-follow-up. Each
  panel reports $\lVert\nabla_\theta L_{\mathrm{dec}}\rVert$ and
  $\lVert\nabla_\theta L_{\mathrm{pred}}\rVert$ on a log scale for
  FDFL at $\mu\in\{0,0.1,0.5,1\}$. Top row: MD knapsack (SPSA)
  shows an order-of-magnitude scale explosion at $\mu=0$ that
  $\mu>0$ suppresses. Bottom row: healthcare (analytic) is flat
  across all $\mu$. Shaded bands are the $5$--$95\%$ seed/lambda
  quantile envelope.}
  \label{fig:gradient-scale-mu-followup}
\end{figure}

\subsection{Takeaway}
The two follow-up experiments isolate two distinct failure modes and
two differently-shaped remedies. The first failure mode---FDFL
divergence under SPSA---is cleanly resolved by adding any positive
fixed prediction weight. The effect is monotone in the SPSA regime
and the choice of $\mu\in\{0.1,0.5,1\}$ is second-order, while in
the analytic regime the anchor is irrelevant. The second failure
mode---PCGrad's scale-domination under SPSA---is only partially
addressed by per-objective normalization: normalization resolves the
directional imbalance on the majority of seeds but leaves a non-trivial
($\approx 11\%$ observed) tail of catastrophic failures, because the
step-size rescaling still transmits gradient-magnitude spikes to the
predictor. On analytic gradients, normalization has an
$\alpha$-dependent effect: a substantial improvement at concave
welfare curvatures and a marginal regression at convex ones. The
practical implication is that a fixed prediction anchor is a safe
default first stabilizer whenever decision gradients are approximate,
while per-objective gradient normalization is best framed as a
configuration option to be enabled selectively rather than as a
drop-in replacement for the original projection.
